\documentclass[10pt,a4paper]{amsart}
\usepackage[margin=19mm]{geometry}
\usepackage{amsmath,amssymb,mathtools}
\usepackage[scr=rsfs,cal=euler]{mathalfa}
\usepackage{xfrac}
\usepackage[unicode,hidelinks,bookmarks=false]{hyperref}
\newcommand{\ie}{i.e.\ }
\newcommand{\cf}{cf.\ }
\newcommand{\resp}{respectively\ }
\newcommand{\Subsection}{\subsection}
\providecommand{\nobreakdash}{\nobreak\hbox{-}}
\setlength{\emergencystretch}{2em}
\multlinegap0pt
\usepackage{amssymb}
\usepackage{mathtools}
\usepackage[shortlabels]{enumitem}
\newtheorem{lemma}{Lemma}
\newtheorem{theorem}{Theorem}
\newcommand{\Fp}{\mathbb{F}_p}
\title{Fibre-Product Stability, Kummer Gluing, and Composita for Universal Koszulity}
\author{Marina Palaisti}
\date{Source: arXiv:2602.09261v4 (CC BY 4.0)}
\begin{document}
\maketitle

\noindent\emph{Source and licence.} Fibre-Product Stability, Kummer Gluing, and Composita for Universal Koszulity. Version arXiv:2602.09261v4 (CC BY 4.0), distributed by arXiv under the Creative Commons Attribution 4.0 International licence, http://creativecommons.org/licenses/by/4.0/, as recorded in the arXiv licence metadata for this exact version and shown on the abstract page https://arxiv.org/abs/2602.09261v4. Full licence notice: \texttt{licenses/arxiv-2602.09261-CC-BY-4.0.txt}.\par\smallskip

\noindent\textit{Editorial scope.} Counted unit: the article's theorem thm:split-raag-lattice (source0.tex lines 1467--1499). The article's complete original proof of it is delivered with it (source0.tex lines 1501--1544, one \texttt{proof} environment of 44 lines). No mathematics is written, repaired or re-derived: the statement and the proof are the article's own lines, in the article's own notation, and no retained line is substituted or re-flowed.  Retained uncounted as the source-local context the proof needs: lines 1347--1362.  Nothing was omitted from any retained span, so the package carries no omission record.  No retained line contains a citation, so the wrapper carries no bibliography and no bibliography entry.  No retained line contains a figure environment, an image-inclusion command or drawing code, so no image is delivered and the package declares no asset dependency.  Editorial formatting only, in addition: an AMS \texttt{amsart} preamble that restates the article's own declarations and macros used by the retained lines, namely the environments \texttt{lemma}, \texttt{theorem} on separate counters, together with the standard packages those lines need.  The retained environments print in this document as Lemma 1, Theorem 1 in order of appearance, whereas the article prints the same items with the numbering of its own full document.  The complete native source of the article as distributed in the arXiv e-print for this exact version is bundled unchanged as \texttt{sources/194307/source0.tex}.
\begin{lemma}
\label{lem:persistence}
Let \(k\) be a field and let \(L/k\) be a finite extension with
\[
k\subseteq L\subseteq k(p).
\]
Then
\[
L(p)=k(p).
\]
Consequently,
\[
G_L(p)=\operatorname{Gal}(k(p)/L)
\]
is an open subgroup of \(G_k(p)\).
\end{lemma}

\begin{theorem}
\label{thm:split-raag-lattice}
Let \(k\) contain \(\mu_p\) and suppose
\[
G_k(p)\cong F_d\times\mathbb Z_p^s,
\qquad d\ge2,\quad s\ge1.
\]
For every finite extension
\[
k\subseteq L\subseteq k(p),
\]
put \(H=G_L(p)\), viewed as an open subgroup of
\(F_d\times\mathbb Z_p^s\), and let
\[
P=\operatorname{pr}_{F_d}(H)\le F_d.
\]
Then \(P\) is open in \(F_d\), and if
\[
q=[F_d:P],
\]
then
\[
G_L(p)\cong F_{\,1+q(d-1)}\times\mathbb Z_p^s.
\]
Consequently
\[
H^\bullet(G_L(p),\Fp)
\cong
A_{\mathcal E_{1+q(d-1)}\ast K_s},
\]
so every field in the finite \(p\)-extension lattice of \(k\) has an
explicit diagonal Stanley--Reisner cohomology presentation.
\end{theorem}

\begin{proof}
By Lemma~\ref{lem:persistence}, \(H=G_L(p)\) is an open subgroup of
\(F_d\times\mathbb Z_p^s\). Write
\[
Z=\mathbb Z_p^s.
\]
Since \(H\) is open, \(H\cap Z\) is open in \(Z\). Every open subgroup
of \(\mathbb Z_p^s\) is abstractly isomorphic to \(\mathbb Z_p^s\), so
\[
H\cap Z\cong\mathbb Z_p^s.
\]
We use this abstract pro-\(p\) group isomorphism in the decomposition below.
Moreover, the image
\[
P=\operatorname{pr}_{F_d}(H)
\]
is open in \(F_d\). Projection gives an exact sequence
\[
1\longrightarrow H\cap Z
\longrightarrow H
\longrightarrow P
\longrightarrow1.
\]
The group \(P\) is an open subgroup of a free pro-\(p\) group and is
therefore free pro-\(p\). Free pro-\(p\) groups are projective in the
category of pro-\(p\) groups, so the displayed sequence splits. Since
\(H\cap Z\) lies in the central factor \(Z\), it is central in \(H\);
therefore the splitting is a direct product:
\[
H\cong P\times(H\cap Z).
\]
The pro-\(p\) Schreier formula gives
\[
P\cong F_{1+q(d-1)},
\qquad q=[F_d:P],
\]
and hence
\[
G_L(p)=H\cong F_{1+q(d-1)}\times\mathbb Z_p^s.
\]
The final statement follows from the standard RAAG cohomology
presentation and the fact that
\(\mathcal E_{1+q(d-1)}\ast K_s\) is diagonal.
\end{proof}

\end{document}
